% !TeX spellcheck = en_US
\documentclass[french]{yLectureNote}

\title{Électromagnétisme 1 PS}
\subtitle{Physique}
\author{Paulhenry Saux}
\date{\today}
\yLanguage{Français}
%lionels.calmels@cemes.fr
\professor{F.Pettinari}
\usepackage{graphicx}%----pour mettre des images
\usepackage[utf8]{inputenc}%---encodage
\usepackage{geometry}%---pour modifier les tailles et mettre a4paper
%\usepackage{awesomebox}%---pour les boites d'exercices, de pbq et de croquis ---d\'esactiv\'e pour les TP de PC
\usepackage{tikz}%---pour deiffner + d\'ependance de chemfig
% \usepackage{tabularx}%---pour dimensionner automatiquement les tableaux avec variable X
\usepackage{awesomebox}%---Pour les boites info, danger et autres
\usepackage{menukeys}%---Pour deiffner les touches de Calculatrice
\usepackage{fancyhdr}%---pour les en-t\^ete personnalis\'ees
\usepackage{blindtext}%---pour les liens
\usepackage{hyperref}%---pour les liens (\`a mettre en dernier)
\usepackage{caption}%---pour la francisation de la l\'egende table vers Tableau
\usepackage{pifont}
\usepackage{array}%---pour les tableaux
\usepackage{lipsum}
\usepackage{yFlatTable}
\usepackage{multicol}
\newcommand{\Lim}[1]{\lim\limits_{\substack{#1}}\:}
\renewcommand{\vec}{\overrightarrow}
\newcommand{\N}[0]{\mathbb{N}}
\newcommand{\dd}{\mathrm{d}}
\newcommand{\norm}[1]{||\vec{#1}||}
\newcommand{\fo}{\psi(\vec{r},t)}
\newcommand{\foe}{\psi(\vec{r},t)\*}
\newcommand{\HH}{\hat{H}}
\newcommand{\hb}{\hbar}
\newcommand{\lap}{\nabla^2}
\newcommand{\lapcc}{\frac{\partial^2 }{\partial x^2}+\frac{\partial^2 }{\partial y^2}+\frac{\partial^2 }{\partial z^2}}
\newcommand{\E}{\vec{E}(M)}
\newcommand{\V}{V(M)}
\newcommand{\er}{\vec{e_{\rho}}}
\newcommand{\ez}{\vec{e_{z}}}
\newcommand{\ep}{\vec{e_{\varphi}}}
\newcommand{\pip}{\pi^+}
\newcommand{\pim}{\pi^-}
\newcommand{\mv}{\mathcal{V}}
\begin{document}
\setcounter{chapter}{3}
\chapter{Distributions continues de charge}
\section{Distribution volumique de charge}
% On note \(\mathcal{V}\) le volume dans lequel est répartie la charge électrique, P un point quelconque dans le volume. On note aussi \(\dd \mv (P)\) est le volume élémentaire situé au point P, \(\rho_C(P)\) la densité volumique de charge au point P et M le point pour lequel on souhaite calculer le champ et potentiel crée par la distribution de charge.
%
%  \(\dd \mv (P)\) contient la charge électrique élémentaire et \(\dd q(P) = \rho_c(P)\dd \mv(P)\). Cette charge quasi-ponctuelle car infiniment petite dans l'espace crée au point M le champ électrique élémentaire \(\dd \E \frac{\dd q(P) \vec{PM}}{4\pi \epsilon_0 \norm{PM}^3}\).
%
%  On en déduit que le champ au point M est donné par \[\E = \iiint_{P\in\mv} \frac{\rho_c(P)\dd \mv(P)\vec{PM}}{4\pi\epsilon_0 \norm{PM}^3}\]
%
%  Cette intégrale triple porte sur les 3 coordonnées permettant de repérer la position de P dans le volume \(\mv\)
%
%  On trouve de la m\^eme façon pour le potentiel :
%  \[\dd \V = \frac{\dd q(P)}{4\pi\epsilon_0 \norm{PM}}\] puis \[V(M) = \iiint_{P\in \mv} \frac{\rho_c \dd \mv(P)}{4\pi\epsilon_0 \norm{PM}}\]
 \begin{theorem}[Champ électrique total  (d. volumique)]
\[\E = \iiint_{P\in\mv} \frac{\rho_c(P)\dd \mv(P)\vec{PM}}{4\pi\epsilon_0 \norm{PM}^3}\]
 \end{theorem}
 \begin{theorem}[Potentiel total  (d. volumique)]
\[V(M) = \iiint_{P\in \mv} \frac{\rho_c \dd \mv(P)}{4\pi\epsilon_0 \norm{PM}}\]
 \end{theorem}
 \section{Distribution surfacique de charge}
On introduit les m\^emes grandeurs que pour la partie précédente, sauf  \(\rho_C(P)\) qui est remplacée par  \(\sigma(P)\) et \(\mv\) par S.

 \begin{theorem}[Champ électrique total (d. surfacique)]
\[\E = \iint_{P\in S} \frac{\sigma(P)\dd S(P)\vec{PM}}{4\pi\epsilon_0 \norm{PM}^3}\]
 \end{theorem}
 \begin{theorem}[Potentiel total  (d. surfacique)]
\[V(M) = \iint_{P\in S} \frac{\sigma \dd S(P)}{4\pi\epsilon_0 \norm{PM}}\]
 \end{theorem}
  \section{Distribution linéique de charge}
On introduit les m\^emes grandeurs que pour la partie précédente, sauf  \(\rho_C(P)\) qui est remplacée par  \(\lambda(P)\) et \(\mv\) par C.

 \begin{theorem}[Champ électrique total (d. linéique)]
\[\E = \int_{P\in C} \frac{\lambda(P)\dd S(P)\vec{PM}}{4\pi\epsilon_0 \norm{PM}^3}\]
 \end{theorem}
 \begin{theorem}[Potentiel total  (d. linéique)]
\[V(M) = \int_{P\in C} \frac{\lambda \dd S(P)}{4\pi\epsilon_0 \norm{PM}}\]
 \end{theorem}
% \section{Exemple : E crée en un point quelconque de son plan médiateur par un segment rectiligne}
%  \subsection{Analyse géométrique}
% La distribution de charge est à symétrie cylindrique autour de l'axe chargé. On repère donc le point par ses coordonnées cylindriques avec un axe Oz confindu avec le segment et O et au milieu du segment.
%
% Si M est dans le plan méditeur, \(z=0\)
%
% Les 2 plans \((m, \er, \vec{e_z}), (m, \er, \ep)\), et \(\E\) est vrai par rapport aux plan, donc il n'y a qu'une composante selon \(\er\).
%
% Il y a une invariance autour de l'axe, donc \(\vec{E} = E_{\rho}(\rho)\) par le principe de Curie.
%
% On introduit le point P appartenant au segment. Il n'est repéré que par \(z_p\).
%
% La longueur élémentaire du segment est \(\dd l\). Pour engendrer le segment chargé, on fait varier \(z_p\) dans es bornes \(z_p+\dd z_p\). \(\dd l(P)\) porte une charge élémentaire \(\dd q = \lambda \dd z_p\).
%
% On en déduit que \[\dd \E = \frac{\lambda}{4\pi\epsilon_0} \frac{\vec{PM}}{\norm{PM}^3}\] avec \(\vec{PM} = \vec{PO} + \vec{OM} = -z_p\vec{e_z}+\rho\er\)
%
% Donc \[\E = \frac{\lambda \dd z_p}{4\pi \epsilon_0}\frac{-z\vec{e_z}+\rho\er}{(\rho^2+z_p^2)^{3/2}}\]
%
% On intègre maintenant entre les bornes pour sommer :
% \[\E = \frac{\lambda}{4\pi\epsilon_0}(\int_{-l/2}^{l/2}\frac{-z_p\vec{e_z}}{(\rho^2+z_p^2)^{3/2}}\dd z_p + \int_{-l/2}^{l/2}\frac{\rho\er}{(\rho^2+z_p^2)^{3/2}})\dd z_p\]
%
% La première inétégrale est nulle car fonction impaire et on a montré qu'il n'y avait pas de composantes selon \(\vec{e_z}\).
%
% Donc \[\E = \frac{\lambda\rho}{4\pi\epsilon_0}\er\int\frac{\dd z_p}{(\rho^2+z_p^2)^{3/2}}\]
%
% On intègre pour trouver \[\E = \frac{\lambda\rho}{4\pi\epsilon_0}\er [\frac{z_p}{\rho^2\sqrt{z_p^2+\rho^2}}]^{l/2}_{-l/2} =  \frac{\lambda}{4\pi\epsilon_0} \frac{z_p}{\rho\sqrt{(l/4)^2+\rho^2}} \er \]
%
% À grande distance, on ne voit pas les effets liés au détail de la distribution de charge finie : on a l'impression de voir une charge ponctuelle \(Q = \lambda l\) au point O.
%
% On retrouve bien le résultat attendu pour une charge ponctuelle.
%
% On peut faire la m\^eme analyse et retrouver un champ électrique crée par un fil infini en faisant tendre \(l\to \infty\) : \[\E = \frac{\lambda}{2\pi \epsilon_0 \rho}\]
% %TODO Refactoriser



\section{Exemple : Champ électrique $\vec{E}$ créé en un point quelconque de son plan médiateur par un segment rectiligne}
\subsection{Analyse Géométrique et Prédiction par Symétrie}
\textbf{Cadre :} Segment de longueur $l$, centré en $O$ sur l'axe $Oz$, chargé uniformément avec une densité linéique $\lambda$. Point d'observation $M$ dans le plan médiateur ($z=0$), repéré par $\rho$.

\begin{enumerate}
    \item \textbf{Coordonnées :} On utilise les coordonnées cylindriques $(\rho, \varphi, z)$. Le point source $P$ est repéré par $z_p \in [-l/2, l/2]$.
    \item \textbf{Invariance et Symétries :}
    \begin{itemize}
        \item \textit{Invariance par Rotation} autour de $Oz$: $\E$ est indépendant de $\varphi$.
        \item \textit{Plans Méridiens} (contenant $Oz$) : plans de symétrie. Le champ $\E$ appartient à ces plans. Donc $E_{\varphi} = 0$.
        \item \textit{Plan Médiateur} ($z=0$) : plan de symétrie. Le champ $\E$ appartient à ce plan. Donc $E_z = 0$.
    \end{itemize}
    \item \textbf{Forme du Champ (Principe de Curie) :} Le champ $\E$ est nécessairement radial :
    $$\E = E_{\rho}(\rho) \er$$
\end{enumerate}

\subsection{Calcul du Champ Électrique par Intégration}

\subsubsection{Champ Élémentaire}
L'élément de charge $\dd q = \lambda \dd z_p$ en $P(z_p)$ crée un champ élémentaire $\dd \E$ en $M(\rho, 0, \varphi)$.
\begin{flalign*}
    \vec{PM} &= \vec{PO} + \vec{OM} = -z_p \ez + \rho \er & \\
    \norm{PM} &= \sqrt{\rho^2 + z_p^2} & \\
    \dd \E &= \frac{1}{4\pi\epsilon_0} \frac{\dd q}{\norm{PM}^3} \vec{PM} = \frac{\lambda \dd z_p}{4\pi\epsilon_0} \frac{\rho \er - z_p \ez}{(\rho^2 + z_p^2)^{3/2}} &
\end{flalign*}

\subsubsection{Intégration et Application de la Symétrie}
Le champ total $\E$ est obtenu par sommation (intégration) de $z_p$ de $-l/2$ à $l/2$ :

\begin{flalign*}
    \E &= \int_{-l/2}^{+l/2} \dd \E & \\
    &= \frac{\lambda}{4\pi\epsilon_0} \left[ \int_{-l/2}^{l/2} \frac{\rho \er}{(\rho^2 + z_p^2)^{3/2}} \dd z_p - \int_{-l/2}^{l/2} \frac{z_p \ez}{(\rho^2 + z_p^2)^{3/2}} \dd z_p \right] & \\
    &\text{Mais } \int_{-l/2}^{l/2} \frac{z_p \ez}{(\rho^2 + z_p^2)^{3/2}} \dd z_p = \vec{0} & \quad \text{(analyse de symétrie)} \\
    &= \frac{\lambda \rho \er}{4\pi\epsilon_0} \int_{-l/2}^{+l/2} \frac{\dd z_p}{(\rho^2 + z_p^2)^{3/2}} &
\end{flalign*}

\subsubsection{Calcul de l'Intégrale}
L'intégrale est calculée en utilisant la primitive $\int \frac{\dd z_p}{(\rho^2 + z_p^2)^{3/2}} = \frac{z_p}{\rho^2 \sqrt{\rho^2 + z_p^2}} + C$.

\begin{flalign*}
    \int_{-l/2}^{+l/2} \frac{\dd z_p}{(\rho^2 + z_p^2)^{3/2}} &= \left[ \frac{z_p}{\rho^2 \sqrt{\rho^2 + z_p^2}} \right]_{-l/2}^{+l/2} & \\
    &= \frac{1}{\rho^2} \left( \frac{l/2}{\sqrt{\rho^2 + (l/2)^2}} - \frac{-l/2}{\sqrt{\rho^2 + (-l/2)^2}} \right) & \\
    &= \frac{1}{\rho^2} \frac{l}{\sqrt{\rho^2 + (l/2)^2}} &
\end{flalign*}

\subsubsection{Résultat Final}
En substituant le résultat de l'intégrale :

\begin{flalign*}
    \E &= \frac{\lambda \rho \er}{4\pi\epsilon_0} \left( \frac{l}{\rho^2 \sqrt{\rho^2 + (l/2)^2}} \right) & \\
    \E &= \frac{\lambda l}{4\pi\epsilon_0 \rho \sqrt{\rho^2 + (l/2)^2}} \er & \quad \text{(Champ créé par le segment fini)}
\end{flalign*}

\subsection{Analyse des Cas Limites}

\subsubsection{Cas du Fil Infini ($l \to \infty$)}
Pour $l \to \infty$, $\sqrt{\rho^2 + (l/2)^2} \approx \sqrt{(l/2)^2} = l/2$.

\begin{flalign*}
    \E_{\text{infini}} &\approx \frac{\lambda l}{4\pi\epsilon_0 \rho (l/2)} \er & \\
    \E_{\text{infini}} &= \frac{\lambda}{2\pi \epsilon_0 \rho} \er & \quad \text{(Résultat pour un fil infini)}
\end{flalign*}

\subsubsection{Cas de la Charge Ponctuelle (Grande Distance $\rho \gg l$)}
Pour $\rho \gg l$, $\sqrt{\rho^2 + (l/2)^2} \approx \sqrt{\rho^2} = \rho$. La charge totale est $Q = \lambda l$.

\begin{flalign*}
    \E_{\text{ponctuel}} &\approx \frac{\lambda l}{4\pi\epsilon_0 \rho (\rho)} \er & \\
    \E_{\text{ponctuel}} &= \frac{1}{4\pi \epsilon_0} \frac{Q}{\rho^2} \er & \quad \text{(Champ d'une charge ponctuelle)}
\end{flalign*}
 \end{document}
